\documentclass[10pt,letterpaper]{article}
\usepackage{cornell-notes}

\title{Clause 28.04.03: Class Template Complex}
\author{}
\date{}
\setDocTitle{Clause 28.04.03: Class Template Complex}
\setDocSubtitle{C++ 2024}
\setDocOwner{C++ 2024 Cornell Notes}
\setDocDate{\today}
\setCornellCollection{C++ 2024 Cornell Notes}
\setCornellUnitType{Clause}
\setCornellUnitNumber{28.04.03}
\setCornellUnitTitle{Class Template Complex}
\hypersetup{pdftitle={Clause 28.04.03: Class Template Complex},pdfsubject={Cornell notes},pdfauthor={}}

\begin{document}
\maketitle
\begin{CornellOverviewBox}{Overview}
\textbf{Classification:} Standard library - Normative\\[2pt]
\textbf{Learning objective:} Understand the rules, terminology, and semantic consequences associated with Class template complex.
\end{CornellOverviewBox}\begin{CornellSummaryBox}{Summary}
{\sffamily\bfseries\color{Accent} Summary}\par\vspace{3pt}Section 28.4.3 focuses on Class template complex. Use the listed grammar productions together with the semantic constraints and disambiguation rules referenced by the section. When applying it, preserve the distinction between mandatory requirements, recommendations, permissions, and behavior deliberately left undefined, unspecified, or implementation-defined.
\end{CornellSummaryBox}

\end{document}
